\section{Resultados IV --- A estrutura intelectual do campo}\label{cap:bibliometria}

Os capítulos anteriores assentam-se no conteúdo \emph{extraído} de cada estudo
pelo modelo de linguagem. Este capítulo triangula aqueles achados por uma lente
independente --- a estrutura de citações do corpus ---, que não depende da
extração assistida por IA. A pergunta é dupla: sobre quais tradições
intelectuais o campo se ergue e como o próprio corpus se divide em grupos
temáticos. Para respondê-la, construíram-se duas redes complementares. Na
\emph{co-citação}, dois trabalhos da literatura de base são ligados quando são
citados em conjunto pelos estudos do corpus, de modo que os agrupamentos revelam
as correntes que fundamentam a área. No \emph{acoplamento bibliográfico}, dois
estudos \emph{do corpus} são ligados quando compartilham referências, de modo
que os agrupamentos revelam sub-literaturas internas. As referências foram
obtidas para 719 dos 817 estudos incluídos (88\%) --- 455 a partir dos próprios
registros da \textit{Web of Science} e 330 via \textit{OpenAlex} ---, unificadas
pelo DOI; estudos e referências sem DOI ficam de fora do pareamento. As
comunidades foram detectadas pelo algoritmo de Louvain. Trata-se de uma análise
\emph{exploratória e descritiva}: como no restante da revisão, o corpus é um
censo e não uma amostra, e as redes ilustram associações, não inferência
populacional.

\subsection{Os fundamentos intelectuais (co-citação)}

A rede de co-citação revela uma base teórica coesa e reconhecível, organizada em
cinco correntes. A primeira, e a mais central, é a tradição da \emph{abordagem
por tarefas} e da polarização do emprego, ancorada na formulação seminal de
\textcite{autorLevyMurnane2003} sobre o conteúdo de tarefas da mudança
tecnológica e desdobrada na evidência de polarização de
\textcite{autorDorn2013}, \textcite{goosManningSalomons2014} e
\textcite{goosManning2007}. A segunda é a corrente de \emph{evidência causal
sobre robôs}, de desenho quase-experimental por região e firma
\parencite{acemogluRestrepo2020robots, graetzMichaels2018, dauth2021}. A
terceira reúne os \emph{fundamentos macroeconômicos e de crescimento} da
automação --- da complementaridade capital-qualificação
\parencite{krusell2000} e do modelo de tarefas de \textcite{zeira1998} ao
declínio da participação do trabalho na renda \parencite{karabarbounisNeiman2014}
e à corrida entre homem e máquina \parencite{acemogluRestrepo2018}. A quarta é a
linha de \emph{risco e prognóstico}, que vai da estimativa de exposição
ocupacional de \textcite{frey2017} e sua revisão por tarefas
\parencite{arntzGregoryZierahn2017} à agenda mais otimista de
\textcite{autor2015} e à evidência inicial sobre IA propriamente dita
\parencite{acemogluVacancies2022, brynjolfssonMitchellRock2018}. A quinta, mais
periférica, é a tradição de \emph{surveys sobre inovação e emprego}
\parencite{vivarelli2014, calvinoVirgillito2017, mondolo2021}. O dado central é
que a literatura recente sobre IA e trabalho não emergiu num vácuo: ela herda uma
espinha teórica madura, da linhagem Autor--Acemoglu, o que confere coerência
conceitual ao campo apesar de sua expansão acelerada.

\subsection{Os grupos temáticos do corpus e o eixo pré/pós}

O acoplamento bibliográfico decompõe os 652 estudos conectados em quatro
sub-literaturas substantivas (Tabela~\ref{tab:acoplamento_clusters}). O maior
grupo (C2, 239 estudos, marcado por \textit{artificial intelligence},
\textit{impact} e \textit{employment}) é nitidamente o da nova literatura sobre
IA: 67\% de seus estudos são posteriores ao ChatGPT. O segundo (C3, 173 estudos,
marcado por \textit{automation}, \textit{job} e \textit{skill}) é o da automação
clássica: 62\% de seus estudos são anteriores ao ChatGPT. Os outros dois ocupam
posições intermediárias --- um voltado à evidência de robôs em nível de firma,
com peso da China (C4), outro à automação, ao crescimento e à desigualdade (C6).

\begin{table}[htbp]
\centering
\caption{Grupos de acoplamento bibliográfico do corpus: período e polarização.}
\label{tab:acoplamento_clusters}
\begin{tabular}{p{6.0cm}ccc}
\toprule
Grupo (termos salientes) & $n$ & \% pós-ChatGPT & \% alta-quali em risco \\
\midrule
C2 --- \textit{artificial intelligence, impact, employment} & 239 & 67\% & 17\% \\
C3 --- \textit{automation, job, market, skill} & 173 & 38\% & 7\% \\
C4 --- \textit{robots, automation, China, firm} & 112 & 57\% & 5\% \\
C6 --- \textit{automation, growth, inequality} & 120 & 60\% & 12\% \\
\bottomrule
\end{tabular}
\par{\footnotesize Acoplamento sobre os 652 estudos conectados (67 isolados
omitidos). A coluna ``\% alta-quali em risco'' é calculada sobre os estudos do
grupo que classificaram a polarização (excluídos os ``n/a''). Análise
exploratória; o corpus é censo, não amostra.}
\end{table}

Esse recorte reproduz, por um caminho inteiramente independente da extração por
IA, o achado central do Capítulo~\ref{cap:comparacao}. O grupo dominado pela
literatura pós-ChatGPT (C2) é justamente o que exibe a \emph{maior} fração de
estudos que situam o risco na alta qualificação (cerca de 17\% dos classificados,
contra apenas 7\% no grupo da automação clássica). Ao mesmo tempo, o risco para a
\emph{baixa} qualificação permanece a categoria modal mesmo em C2 (47\% dos
classificados). Ou seja, a própria geografia de citações do campo desenha um
deslocamento parcial do risco em direção às ocupações cognitivas qualificadas,
sem que esse deslocamento se torne dominante --- exatamente a leitura de
\emph{ruptura parcial} sustentada pela análise de conteúdo. A convergência de dois
métodos distintos sobre a mesma conclusão reforça a confiança no diagnóstico.
Cabe a ressalva, contudo, de que o acoplamento é dominado por um grande
aglomerado geral e de que toda a leitura é exploratória: ela sugere uma
estrutura, não a testa.
